% TRANSLATION UNIT: translate this file only.
% ===========================================================================
%  08 — Modelo de datos
%  The \section title is the canonical (Spanish) chapter name; the body prose
%  is written in ENGLISH and translated to Spanish later.
% ===========================================================================

\section{Modelo de datos}

% DONE: The chapter argues the invariant (database-per-service, no cross-service
% reads) and the four hard shapes (exact decimals, the formula AST in JSONB, the
% article body as a block document, daterange indicators with EXCLUDE gist), plus
% identity/lifecycle, content-addressed images, the deliberately-absent schema
% features, and one entity table per owning service.

% Claims inventory. Every claim in the prose below is checkable in the
% repository; the migrations are the authority wherever a spec and SQL disagree.
% EVIDENCE: database-per-service, no cross-service access → .specify/memory/constitution.md:89
% EVIDENCE: money/tax/rates as arbitrary-precision decimals, float forbidden → .specify/memory/constitution.md:177
% EVIDENCE: versioned paired migrations, no synchronize:true → .specify/memory/constitution.md:282
% EVIDENCE: multi-table writes only through execTx(ctx, fn) → services/users/internal/storer/storer_postgres.go:50; services/learning/src/common/tx.ts:42
% EVIDENCE: simulations store canonical decimal strings, jsonb_has_no_numbers guard → services/simulator/migrations/20260727120000_init_simulator.up.sql:22-68
% EVIDENCE: pot exact for integer exponents, potd approximate for decimal, kept separate → services/simulator/src/domain/formula/functions.rs:190-193
% EVIDENCE: calculators + calculator_definitions; AST persisted as JSONB → services/simulator/migrations/20260902120000_calculators.up.sql:16-63
% EVIDENCE: published_version points at an existing definition, deferrable FK → services/simulator/migrations/20260902130000_calculator_published_version.up.sql:65-74
% EVIDENCE: ten seed definitions (code is the authority) → services/simulator/src/domain/seeds/mod.rs:144-155
% EVIDENCE: financial_indicators NUMERIC(20,6) with EXCLUDE USING gist → services/simulator/migrations/20260902121500_financial_indicators.up.sql:19-44
% EVIDENCE: simulations snapshot calculator_id/version and indicators_snapshot → services/simulator/migrations/20260902123000_simulations_calculator_ref.up.sql:14-45
% EVIDENCE: article body as JSONB block document, closed vocabulary validated server-side → services/learning/migrations/20260902111500_article_versions_body_doc.up.sql; services/learning/src/articles/body-doc.ts:26-57
% EVIDENCE: article_images BYTEA, SHA-256 id, 2 MB cap, STORAGE EXTERNAL → services/learning/migrations/20260902113000_article_images.up.sql:44-58
% EVIDENCE: four roles usuario_final/editor/coordinador_editorial/administrador → services/users/migrations/20260727120000_init_users.up.sql:45; services/users/migrations/20260903090000_roles_administrador.up.sql
% EVIDENCE: audit_log append-only via REVOKE plus trigger → services/audit/migrations/20260727120000_init_audit.up.sql:58-90
% EVIDENCE: pending_deletion and purge grace specified but not yet migrated → specs/002-calculator-builder-content-admin/data-model.md:315-327 vs services/users/migrations/ (no such migration)

A data model is not a list of tables; it is the set of shapes a system is
willing to commit to, and every shape in FintCart traces back to one decision or
one prohibition. The prohibition comes first. The system is a set of eight
backend services and one Angular SPA, and each stateful backend owns its own
PostgreSQL 16 database that no other service may read or write
(\code{.specify/memory/constitution.md:89}, Principle~III). Seven databases
exist; the API Gateway and the frontend hold none. The consequence that governs
the whole chapter is that an entity is never normalized across a service
boundary, even when it is duplicated: the same e-mail address lives in
\code{auth_db.credentials} and in \code{users_db.profiles}, the same person's
identifier appears in every history table, and in no case is there a foreign key
joining two databases. Cross-service references are opaque UUIDs resolved over
gRPC, and the two profiles agree because the saga that creates the account
writes both, not because either table knows the other exists.

That duplication is bought on purpose, and it is worth naming the price. The
system pays for the isolation with redundant storage and with eventual
consistency: when a quiz is graded, the learning service stores the attempt and
a later saga step updates the users service's best score, so the two services
can disagree for the length of that step. The design accepts that window because
the alternative --- one shared schema that any service can join --- is precisely
the coupling that Principle~III exists to prevent. Duplication is the tax on
sovereignty, and the schemas below are shaped to make the duplicated copies
self-consistent on their own, with the cross-service agreement delegated to
sagas and events rather than to constraints.

\subsection{The four hard shapes}

Four column shapes carry the non-negotiable parts of the model, and each exists
because of a specific failure it rules out.

\textbf{Money and rates as exact decimals.} The single thing the constitution
forbids outright is a binary floating-point number standing for money, a rate,
or a simulation result (\code{.specify/memory/constitution.md:177},
Principle~VIII). The rule names its representation per technology:
\code{NUMERIC} columns in PostgreSQL, \code{shopspring/decimal} in Go,
\code{rust_decimal} in Rust, \code{decimal.js} or \code{big.js} in TypeScript,
and a canonical decimal \code{string} in Protocol Buffers and JSON. The
simulator service enforces the rule at the schema level as well as in code: the
\code{simulations} table stores its \code{inputs} and \code{result} documents as
\code{JSONB}, and a helper function, \code{jsonb_has_no_numbers}, walks the
document to any depth and rejects a raw JSON number anywhere in it
(\code{services/simulator/migrations/20260727120000_init_simulator.up.sql:22-68}).
The reason is subtle and worth stating: a JSON number would be stored as a
PostgreSQL \code{numeric}, but deserializing it into Rust would pass through an
\code{f64} in most deserializers, and that single hop reintroduces the binary
float the rule forbids. The cost of the rule is permanent: every boundary needs
a conversion, and the mapping layer is where it must happen and only there
(Principle~IX). A decimal string, a domain decimal, and a database \code{numeric}
are three distinct types, and the conversion between them is the highest-risk
operation in the system, so it is confined to one module per boundary instead of
scattered through the code.

\textbf{The formula as a persisted AST.} A calculator in the simulator service
is a row, not code. The \code{calculators} table holds the identity --- owner,
curation state, the version in force --- while \code{calculator_definitions}
holds one immutable row per version, with the parsed formula tree stored as
\code{JSONB} next to the source text the author typed
(\code{services/simulator/migrations/20260902120000_calculators.up.sql:16-63}).
Editing never updates a definition; it inserts a new version, which is what lets
a simulation cite the exact version that produced it and what lets a published
calculator keep serving its approved version while the author drafts the next.
Storing the tree rather than the text is what buys user-defined calculators
without deploying: a formula arrives as text, is parsed and validated at save
time, and from then on only the tree is walked, so execution never parses text
and never runs code. What the shape demands is a bounded evaluator: a definition
may declare at most 20 inputs and 10 outputs, and each formula at most 64 nodes
and depth 16, checked by the parser before persistence --- limits the schema
deliberately does not re-encode, because walking a tree to count nodes inside a
\code{CHECK} would be unreadable and would drift from the parser that is the
real gate. Two functions are kept apart on purpose: \code{pot} raises to an
integer exponent by repeated multiplication and is exact, while \code{potd}
raises to a decimal exponent through \code{exp}/\code{ln} and is approximate
(\code{services/simulator/src/domain/formula/functions.rs:190-193}). Merging
them would silently lose the exactness of Principle~VIII, so they are two names
and the author chooses. One count must be corrected here: the simulator ships
ten seed definitions --- \code{ahorro}, \code{credito}, \code{presupuesto},
\code{inversion}, \code{ea_a_mv}, \code{mv_a_ea}, \code{gmf},
\code{fondo_emergencia}, \code{meta_ahorro} and \code{cdt_dias}
(\code{services/simulator/src/domain/seeds/mod.rs:144-155}) --- while the 002
data model still says seven (\code{specs/002-calculator-builder-content-admin/data-model.md:176,300})
and even the migration's own comment repeats that stale count
(\code{20260902120000_calculators.up.sql:18}). The code is what runs, so the
spec is one count behind the system it describes; the lag is reported here
rather than silently reconciled.

\textbf{The article body as a block document.} The learning service replaced the
plain-text \code{article_versions.body} column with \code{body_doc}, a
\code{JSONB} tree of typed blocks with a closed vocabulary --- \code{doc},
\code{parrafo}, \code{encabezado}, \code{lista}, \code{item_lista},
\code{texto}, \code{imagen} and \code{calculadora} --- and a closed set of marks,
\code{negrita}, \code{cursiva} and \code{enlace}
(\code{services/learning/src/articles/body-doc.ts:26-57}). The whole document is
validated positively in the server at save time: the validator enumerates what
is admitted and rejects everything else, links included only if their scheme is
\code{http}, \code{https} or \code{mailto}. The rejected alternative was HTML,
which would have forced the frontend to render with \code{innerHTML} and
\code{bypassSecurityTrust*}, opening an injection surface that then has to be
sanitized forever. A tree of typed blocks has no such surface: the frontend
renders each node with a component and never interpolates markup, so safety is
by construction rather than by sanitizer. The migration that installed the shape
wrapped every existing paragraph in a document and only dropped \code{body}
after a guard proved the converted text equaled the original
(\code{services/learning/migrations/20260902150000_drop_article_versions_body.up.sql}).
The database enforces only the root --- \code{body_doc} must be an object whose
\code{tipo} is \code{doc} --- and leaves the full vocabulary to the validator,
because a complete schema written in SQL would be unreadable and would drift from
the validator that is the real gate.

\textbf{Indicators with a validity range.} A yearly indicator such as \code{@UVT}
does not hold one value; it holds a value for an interval, and the schema makes
the interval part of the row. \code{financial_indicators} stores each value as
\code{NUMERIC(20,6)} together with a half-open \code{DATERANGE} validity, and the
constraint \code{EXCLUDE USING gist (name WITH =, validity WITH &&)} makes the
database itself reject two overlapping values for the same indicator
(\code{services/simulator/migrations/20260902121500_financial_indicators.up.sql:19-44}).
The exclusion lives in the base rather than in the application because two
administrators loading the same year's value concurrently is a real case, and a
read-then-write check in code does not see the race. The half-open convention is
what makes \enquote{which value rules on day D} have a single answer. Because
values change over time, every simulation stores the snapshot it used ---
\code{indicators_snapshot}, a map of name to value written as decimal strings
and guarded by the same no-JSON-numbers check as \code{inputs} and \code{result}
--- instead of joining the indicator table later
(\code{services/simulator/migrations/20260902123000_simulations_calculator_ref.up.sql:14-45}).
A historical simulation therefore stays explicable even after an indicator is
revised, which a join against the current value could never guarantee.

\subsection{Identity and lifecycle}

The identity shapes live in two services and are joined only by the opaque
identifier. \code{auth_db.credentials} owns the account: the \code{id} is the
same UUID as \code{users_db.profiles.id} and as the JWT \code{sub} claim, the
password is an Argon2id hash, and the verification token for e-mail confirmation
is stored only as a SHA-256 hash so that a dump of the table cannot activate a
pending account. \code{users_db.profiles} owns the profile and carries the four
authorization roles through \code{roles_assignment}: \code{usuario_final},
\code{editor}, \code{coordinador_editorial} and \code{administrador}
(\code{services/users/migrations/20260903090000_roles_administrador.up.sql}).
The administrator role is independent of the editorial coordinator; it does not
inherit the coordinator's powers, and the reverse is equally true.

The lifecycle of an account is the place where the specification and the
migrations disagree, and the disagreement is stated here because the SQL is what
runs. The 002 data model specifies a \code{pending_deletion} account status with
a 30-day grace period, two new columns (\code{purge_due_at},
\code{purge_requested_by}), and a widened partial unique index
\code{profiles_email_reserved_uniq} so the e-mail stays reserved during the
grace (\code{specs/002-calculator-builder-content-admin/data-model.md:315-327}).
No users migration implements any of that: the live schema still constrains
\code{account_status} to \code{active} and \code{anonymized}, and the e-mail
uniqueness index still covers only \code{active} accounts
(\code{services/users/migrations/20260727120000_init_users.up.sql:26-32}).
The grace-period shape is therefore specified but not yet migrated, and the
frontend profile component confirms it in a comment. One part of that
specification is already decided and must not be reintroduced by the eventual
migration: the system will not keep a hash of the deleted original e-mail,
because a hash of an address is a membership oracle over an enumerable,
low-entropy space, and anyone with the table could test whether a specific
person ever held an account --- which is exactly the re-identification the
anonymization destroys.

\subsection{Images addressed by content}

Article images live in the database as \code{BYTEA}, with the SHA-256 of the
bytes as the primary key, a 2\,MB ceiling, and the column marked
\code{STORAGE EXTERNAL} so the payload is not compressed into the row
(\code{services/learning/migrations/20260902113000_article_images.up.sql:44-58}).
Addressing by content makes two failures impossible at once: the same image
uploaded twice resolves to the same row, so bytes are never duplicated, and a
file can never be swapped for a different file while keeping its identifier,
because the identifier is the content. The row is immutable by construction ---
there is no \code{updated_at} and no update path --- and the alternative text and
caption live in the \code{imagen} node of the document rather than on the row,
because they belong to a particular version's use of the image, which is what
lets two versions share one file with different captions. The cost is equally
deliberate: there is no CDN and no partial read, a listing query must avoid
pulling the \code{BYTEA} column, and the 2\,MB cap is a schema-enforced promise,
checked against the real byte length with \code{byte_size = octet_length(bytes)}
so that a declared size cannot lie about the content it describes.

\subsection{What the schemas deliberately do not do}

The negative space of the model is as deliberate as the columns. There are no
foreign keys across services: every cross-service reference is an opaque UUID
without a constraint, because the constraint would be the first thread of a
shared schema. There are no triggers carrying business logic; the only triggers
in the system exist to make \code{audit_log} refuse \code{UPDATE} and
\code{DELETE}, enforcing append-only immutability in the base itself
(\code{services/audit/migrations/20260727120000_init_audit.up.sql:58-90}).
There is no \code{synchronize: true} anywhere: every schema change is a paired
versioned migration applied with \code{golang-migrate}, so no ORM silently
rewrites a database on startup (Principle~XI). And multi-table writes run only
through the \code{execTx(ctx, fn)} helper --- in Go
(\code{services/users/internal/storer/storer_postgres.go:50}) and in TypeScript
(\code{services/learning/src/common/tx.ts:42}) --- which concentrates
\code{begin}/\code{rollback}/\code{commit} in one place instead of repeating it
per method. Each of these absences is a decision with a name in the
constitution; together they are what keeps the seven databases seven separate
systems rather than one.

\subsection{Entities per service}

The tables that follow consolidate the model: one table per owning service,
listing the entities and the invariant each carries. The Gateway and the
frontend are omitted because they hold no database. Each table is introduced
briefly; the invariants are the constraints written into the migrations, not a
summary of the application layer.

\begin{table}[htbp]
  \centering
  \caption{Authentication entities (\texttt{auth\_db}) and their invariants.}
  \label{tab:auth}
  \begin{tabular}{@{}l l p{8.2cm}@{}}
    \hline
    \textbf{Entity} & \textbf{Holds} & \textbf{Invariant it carries} \\
    \hline
    \code{credentials} & identity, Argon2id hash & e-mail unique; \code{login_status} in \{\code{pending_verification}, \code{active}, \code{anonymized}\}; verification token stored only as SHA-256 \\
    \code{oauth_clients} & OAuth2 client registry & public SPA has no secret, confidential M2M must; grant types restricted to the two constitutional flows \\
    \code{authorization_codes} & single-use PKCE codes & \code{S256} only; lifetime at most 60 seconds \\
    Redis (not a table) & JWT blacklist, refresh tokens & confined to Principle~IV purposes \\
    \hline
  \end{tabular}
\end{table}

The authentication service owns identity and credentials, and the invariants
above are the trust boundary: an e-mail must be unique, a public client cannot
hold a secret, and a short-lived code cannot outlive its purpose.

\begin{table}[htbp]
  \centering
  \caption{Users entities (\texttt{users\_db}) and their invariants.}
  \label{tab:users}
  \begin{tabular}{@{}l l p{7.6cm}@{}}
    \hline
    \textbf{Entity} & \textbf{Holds} & \textbf{Invariant it carries} \\
    \hline
    \code{profiles} & profile, e-mail copy & \code{id} equals \code{credentials.id} with no cross-service key; \code{account_status} in \{\code{active}, \code{anonymized}\} today \\
    \code{roles_assignment} & role grants & role in the four-value vocabulary \\
    \code{preferences} & locale, channels & one row per user \\
    \code{progress} & points & non-negative; equals the sum of best scores per distinct quiz \\
    \code{quiz_best_score} & best per quiz & updated only when the new score exceeds the stored one \\
    \code{article_views} & view counters & one row per user and article \\
    \code{inapp_notifications} & inbox & \code{read_at} set if and only if \code{read_state} is \code{read} \\
    \hline
  \end{tabular}
\end{table}

The users service owns the profile and every derived reading of a person's
activity. Its most instructive invariant is \code{quiz_best_score}: it is updated
monotonically, which is what makes the saga step that applies a quiz score
idempotent and removes the need for a destructive compensation.

\begin{table}[htbp]
  \centering
  \caption{Learning entities (\texttt{learning\_db}) and their invariants.}
  \label{tab:learning}
  \begin{tabular}{@{}l l p{7.6cm}@{}}
    \hline
    \textbf{Entity} & \textbf{Holds} & \textbf{Invariant it carries} \\
    \hline
    \code{articles} & article identity, category & category references the catalog, never free text \\
    \code{article_versions} & immutable bodies & approver differs from author; \code{body_doc} root is a \code{doc} \\
    \code{categories} & catalog & unique slug; never physically deleted, only deactivated \\
    \code{article_images} & image bytes & primary key is the SHA-256; at most 2\,MB; \code{STORAGE EXTERNAL} \\
    \code{quizzes} & the instrument & \code{pass_threshold} in 0--100; \code{questions_to_serve} positive \\
    \code{questions} & question bank & correct key exists among the options; weight positive \\
    \code{quiz_sessions} & served random set & expires 60 minutes after creation; served order kept \\
    \code{quiz_attempts} & graded submissions & \code{score} in 0--100; \code{served_snapshot} not null; session key \code{SET NULL} \\
    \code{article_stats} & aggregates & counts non-negative \\
    \hline
  \end{tabular}
\end{table}

The learning service owns content and assessment, and its invariants are the
editorial rules made mechanical: nobody approves their own work, a quiz cannot
serve a question whose correct answer is not among the options, and an attempt
keeps its own copy of what was served so the history stays reconstructible after
its session is purged. One migration-level correction is worth recording: the
002 data model prescribed rescaling existing \code{pass_threshold} and
\code{score} from an absolute weighted sum to a percentage
(\code{specs/002-calculator-builder-content-admin/data-model.md:117,158}), but
the actual migrations only tightened the range to 0--100, because the 001 code
had already stored both as percentages and reapplying the formula would have
doubled them (\code{services/learning/migrations/20260902103000_quizzes_questions_to_serve.up.sql}).

\begin{table}[htbp]
  \centering
  \caption{Simulator entities (\texttt{simulator\_db}) and their invariants.}
  \label{tab:simulator}
  \begin{tabular}{@{}l l p{7.6cm}@{}}
    \hline
    \textbf{Entity} & \textbf{Holds} & \textbf{Invariant it carries} \\
    \hline
    \code{calculators} & identity, curation state & published implies an approver and a published version; approver differs from owner; built-in implies no owner \\
    \code{calculator_definitions} & immutable version rows & composite key \code{(calculator_id, version)}; 1--20 inputs, 1--10 outputs; AST bounded by the parser \\
    \code{financial_indicators} & yearly values & \code{NUMERIC(20,6)}; half-open \code{DATERANGE}; no two overlapping values per name \\
    \code{simulations} & history, snapshot & \code{calculator_id} plus exact version; \code{indicators_snapshot}; no JSON number at any depth \\
    \hline
  \end{tabular}
\end{table}

The simulator service owns the formula engine and the history it produces. Its
sharpest invariant is the deferrable foreign key that ties
\code{calculators.published_version} to an existing definition row, which closes
the only path by which a published calculator could point at a version nobody
approved.

\begin{table}[htbp]
  \centering
  \caption{Notification entities (\texttt{notification\_db}) and their invariants.}
  \label{tab:notification}
  \begin{tabular}{@{}l l p{7.6cm}@{}}
    \hline
    \textbf{Entity} & \textbf{Holds} & \textbf{Invariant it carries} \\
    \hline
    \code{notification_events_queue} & pending e-mail & \code{event_id} unique; template in a closed set \\
    \code{notification_states} & outcome that survives dequeue & \code{failed} requires an error; \code{sent} requires none \\
    \hline
  \end{tabular}
\end{table}

The notification service keeps a persistent queue with state: the queue holds
what is pending and the state table survives the dequeue, so the outcome of each
delivery attempt remains auditable after the event leaves the queue.

\begin{table}[htbp]
  \centering
  \caption{Audit entities (\texttt{audit\_db}) and their invariants.}
  \label{tab:audit}
  \begin{tabular}{@{}l l p{7.6cm}@{}}
    \hline
    \textbf{Entity} & \textbf{Holds} & \textbf{Invariant it carries} \\
    \hline
    \code{audit_log} & immutable record & append-only enforced by \code{REVOKE} and a trigger; \code{actor_ref} opaque; partitioned by year \\
    \hline
  \end{tabular}
\end{table}

The audit service owns the one table that must never change: immutability is
enforced in the base, not in the application, and the actor reference is an
opaque identifier that survives the anonymization of the person it names.

\begin{table}[htbp]
  \centering
  \caption{Orchestrator entities (\texttt{orchestrator\_db}) and their invariants.}
  \label{tab:orchestrator}
  \begin{tabular}{@{}l l p{7.6cm}@{}}
    \hline
    \textbf{Entity} & \textbf{Holds} & \textbf{Invariant it carries} \\
    \hline
    \code{saga_state} & saga sequencing & \code{saga_type} and \code{status} in closed vocabularies \\
    \code{event_outbox} & reliable publication & \code{event_type} in a closed set; written in the same transaction as the advance \\
    \hline
  \end{tabular}
\end{table}

The orchestrator holds no domain data --- only the state of sequences and the
outbox that guarantees an event is not lost between a commit and a publish. Its
two closed vocabularies grow only by migration, which is the friction that keeps
an event from being published to a consumer that does not exist.
